% preamble-exam-zh.tex — 共享 preamble
% 用于所有 exam-zh 模板
%
% 样式策略：
%   屏幕 → 语义彩色（蓝=关键想法，红=易错，绿=路线，灰=提示/教师备注）
%   打印 → 自动降级为黑白（通过 print mode 或 monochrome 选项）

\documentclass{exam-zh}

% ---------- 基础配置 ----------
\examsetup{
  page/size = a4paper,
  page/show-head = false,
  page/show-foot = false,
  sealline/show = false,
  font = times,
  math-font = xits,
}

\usepackage{tcolorbox}
\usepackage{needspace}
\tcbuselibrary{skins,breakable}

% ---------- 打印/屏幕颜色切换 ----------
% 用法：\documentclass[monochrome]{exam-zh} 即可切换为纯黑白
% 注意：不要使用 \if@xxx 放在 makeatother 外部；这里改成普通条件名。
\newif\ifeduprintcolor
\eduprintcolortrue

\makeatletter
\@ifclasswith{exam-zh}{monochrome}{\eduprintcolorfalse}{}
\makeatother

% ---------- 语义颜色定义 ----------
% 屏幕：有色彩区分；打印：降级为灰度
\ifeduprintcolor
  % --- 屏幕彩色模式 ---
  \definecolor{edu-blue}{RGB}{47, 82, 122}       % 关键想法
  \definecolor{edu-blue-bg}{RGB}{243, 246, 252}
  \definecolor{edu-red}{RGB}{180, 40, 40}         % 易错提醒
  \definecolor{edu-red-bg}{RGB}{255, 245, 240}
  \definecolor{edu-green}{RGB}{34, 100, 60}       % 路线图
  \definecolor{edu-green-bg}{RGB}{242, 250, 245}
  \definecolor{edu-orange}{RGB}{160, 90, 0}       % 训练目标
  \definecolor{edu-orange-bg}{RGB}{255, 248, 235}
  \definecolor{edu-gray}{RGB}{100, 100, 100}      % 提示/教师备注
  \definecolor{edu-gray-bg}{RGB}{248, 248, 248}
  \definecolor{edu-purple}{RGB}{90, 50, 120}      % 思考题
  \definecolor{edu-purple-bg}{RGB}{248, 244, 252}
\else
  % --- 打印黑白模式 ---
  \definecolor{edu-blue}{gray}{0.15}
  \definecolor{edu-blue-bg}{gray}{0.96}
  \definecolor{edu-red}{gray}{0.25}
  \definecolor{edu-red-bg}{gray}{0.97}
  \definecolor{edu-green}{gray}{0.20}
  \definecolor{edu-green-bg}{gray}{0.97}
  \definecolor{edu-orange}{gray}{0.25}
  \definecolor{edu-orange-bg}{gray}{0.98}
  \definecolor{edu-gray}{gray}{0.40}
  \definecolor{edu-gray-bg}{gray}{0.97}
  \definecolor{edu-purple}{gray}{0.30}
  \definecolor{edu-purple-bg}{gray}{0.98}
\fi

% ---------- 自定义教学环境 ----------
% 共性：enhanced + breakable（跨页不断裂）

% 原题卡片 — 强边框，突出显示
\newtcolorbox{problemcard}{
  enhanced, breakable,
  colback=edu-blue-bg,
  colframe=edu-blue,
  boxrule=1.2pt,
  arc=0pt,
  left=3mm, right=3mm, top=3mm, bottom=3mm,
}

% 解题路线图 — 左边线 + 浅底
\newtcolorbox{routebox}{
  enhanced, breakable,
  colback=edu-green-bg,
  colframe=edu-green!30,
  boxrule=0pt,
  borderline west={2.5pt}{0pt}{edu-green},
  left=3mm, right=3mm, top=3mm, bottom=3mm,
}

% 关键想法 — 蓝色左边线
\newtcolorbox{keyideabox}{
  enhanced, breakable,
  colback=edu-blue-bg,
  colframe=edu-blue!30,
  boxrule=0pt,
  borderline west={2.5pt}{0pt}{edu-blue},
  left=3mm, right=3mm, top=3mm, bottom=3mm,
}

% 易错提醒 — 红色左边线
\newtcolorbox{mistakebox}{
  enhanced, breakable,
  colback=edu-red-bg,
  colframe=edu-red!20,
  boxrule=0pt,
  borderline west={3pt}{0pt}{edu-red},
  left=3mm, right=3mm, top=2.5mm, bottom=2.5mm,
}

% 提示 — 灰色虚线左边线
\newtcolorbox{hintbox}{
  enhanced, breakable,
  colback=edu-gray-bg,
  colframe=edu-gray!20,
  boxrule=0pt,
  borderline west={2pt}{0pt}{edu-gray, dashed},
  left=3mm, right=3mm, top=2mm, bottom=2mm,
}

% 思考题 — 紫色虚线左边线
\newtcolorbox{thinkbox}{
  enhanced, breakable,
  colback=edu-purple-bg,
  colframe=edu-purple!20,
  boxrule=0pt,
  borderline west={2pt}{0pt}{edu-purple, dashed},
  left=2.5mm, right=2.5mm, top=2.5mm, bottom=2.5mm,
}

% 教师备注 — 灰色左边线，小字
\newtcolorbox{teachernote}{
  enhanced, breakable,
  colback=edu-gray-bg,
  colframe=edu-gray!15,
  boxrule=0pt,
  borderline west={2pt}{0pt}{edu-gray!60},
  left=2.5mm, right=2.5mm, top=2.5mm, bottom=2.5mm,
  fontupper=\small,
  colupper=black!60,
}

% 训练目标 — 橙色左边线，小字
\newtcolorbox{traininggoal}{
  enhanced, breakable,
  colback=edu-orange-bg,
  colframe=edu-orange!15,
  boxrule=0pt,
  borderline west={1.5pt}{0pt}{edu-orange},
  left=2mm, right=2mm, top=1mm, bottom=1mm,
  fontupper=\small,
}

% 答题区 — 无边框，纯留白
\newcommand{\answerarea}[1][40mm]{%
  \vspace{2mm}%
  \begin{tcolorbox}[
    enhanced, breakable,
    colback=white,
    colframe=white,
    boxrule=0pt,
    height=#1,
    valign=top,
    bottom=0mm,
    after skip=0mm,
  ]
  \end{tcolorbox}%
}

% ---------- 字体设置 ----------
% exam-zh (ctexbook) already sets CJK fonts via ctex.
% Only override if specific fonts are needed.
% macOS: Songti SC, STHeiti, STFangsong
% Linux: SimSun, SimHei, FangSong

% ---------- 页面设置 ----------
% exam-zh (ctexbook) already loads geometry.
% Override margins if needed:
% \geometry{top=16mm, bottom=16mm, left=14mm, right=14mm}

\examsetup{
  question/show-answer = true,
  fillin/show-answer = true,
  solution/show-solution = show-stay,
}

% ---------- 教师版解答样式 ----------
\newtcolorbox{solutionblock}{
  enhanced, breakable,
  colback=edu-blue-bg,
  colframe=edu-blue!25,
  boxrule=0pt,
  borderline west={2pt}{0pt}{edu-blue!50},
  arc=0pt,
  left=3mm, right=3mm, top=2mm, bottom=2mm,
  before skip=2mm,
  after skip=3mm,
  fontupper=\small,
}

% 解答步骤编号
\newcounter{solstep}
\newcommand{\solstep}[2]{%
  \stepcounter{solstep}%
  \par\medskip
  \noindent\hangindent=6mm
  {\color{edu-blue}\bfseries\textcircled{\scriptsize\thesolstep}\enspace #1}\enspace #2\par
}

\begin{document}

\section*{通过交点、面积、距离等条件求点坐标（一） · 练习卷（教师版）}

\section*{一、选择题（每题 12 分，共 48 分）}


\begin{keyideabox}
  \textbf{核心思想：直线过定点与交点}
  1. 直线过点 $(x_0, y_0)$ $\Leftrightarrow$ 点坐标满足解析式；\\ 2. 两直线交点坐标，是联立两直线解析式方程组的解。
\end{keyideabox}



\begin{question}[points=12]
  直线 $y = 2x + b$ 经过点 $(3, 8)$，则 $b$ 的值为
  \begin{choices}
    \item 2
    \item -2
    \item 14
    \item 22
  \end{choices}
  \paren[A]
  \begin{solutionblock}
    将 $(3, 8)$ 代入 $y = 2x + b$ 得：$8 = 2 \times 3 + b$，解得 $b = 2$。
  \end{solutionblock}
\end{question}



\begin{question}[points=12]
  直线 $y = 2x + 1$ 与直线 $y = -x + 4$ 的交点坐标为
  \begin{choices}
    \item $(1, 3)$
    \item $(3, 1)$
    \item $(2, 5)$
    \item $(5, 2)$
  \end{choices}
  \paren[A]
  \begin{solutionblock}
    联立得 $2x + 1 = -x + 4 \implies 3x = 3 \implies x = 1$。代入得 $y = 3$。所以交点为 $(1, 3)$。
  \end{solutionblock}
\end{question}



\begin{question}[points=12]
  直线 $y = -3x + 6$ 与 $x$ 轴交于点 $A$，过原点 $O$ 作该直线的平行线与 $y$ 轴交于点 $B$，则 $B$ 的坐标为
  \begin{choices}
    \item $(2, 0)$
    \item $(0, 6)$
    \item $(0, 0)$
    \item $(0, -6)$
  \end{choices}
  \paren[C]
  \begin{solutionblock}
    平行线斜率相同，故平行线为 $y = -3x$。该线过原点，与 $y$ 轴交点即为 $(0,0)$。
  \end{solutionblock}
\end{question}



\begin{question}[points=12]
  点 $A(0, 3)$，点 $B$ 在第一象限内，且线段 $AB \perp y$ 轴。若 $AB = 5$，则点 $B$ 的坐标为
  \begin{choices}
    \item $(5, 3)$
    \item $(3, 5)$
    \item $(5, 8)$
    \item $(3, 0)$
  \end{choices}
  \paren[A]
  \begin{solutionblock}
    $AB \perp y$ 轴说明 $B$ 纵坐标与 $A$ 相同为 $3$；$B$ 在第一象限且距离为 $5$，故横坐标为 $5$。坐标为 $(5, 3)$。
  \end{solutionblock}
\end{question}


\section*{二、填空题（每题 12 分，共 12 分）}


\begin{question}[points=12]
  已知平面直角坐标系中两点 $A(1, 2)$ 和 $B(4, 6)$，则线段 $AB$ 的长度为 \fillin。
  \paren[5]
  \begin{solutionblock}
    根据两点间距离公式：$AB = \sqrt{(4-1)^2 + (6-2)^2} = \sqrt{3^2 + 4^2} = 5$。
  \end{solutionblock}
\end{question}


\section*{三、解答题（共 40 分）}

\needspace{8\baselineskip}

\begin{problem}[points=20]
  直线 $y = -\dfrac{4}{3}x + 4$ 与 $x$ 轴、$y$ 轴分别交于点 $A$、$B$。

(1) 求点 $A$ 和点 $B$ 的坐标；
(2) 若点 $P$ 在 $x$ 轴上，且 $S_{\triangle POB} = 6$，求点 $P$ 的坐标。
  \answerarea[50mm]
  \setcounter{solstep}{0}
  \begin{solutionblock}
    \textbf{答案：}(1) $A(3, 0)$，$B(0, 4)$；(2) $P(3, 0)$ 或 $P(-3, 0)$\par\medskip
    \solstep{求 A、B 坐标}{令 $y = 0 \implies -\frac{4}{3}x + 4 = 0 \implies x = 3$，故 $A(3, 0)$。\\ 令 $x = 0 \implies y = 4$，故 $B(0, 4)$。}
    \solstep{列出面积方程}{点 $P(x_p, 0)$ 在 $x$ 轴上，$OB = 4$ 在 $y$ 轴上，高为 $|x_p|$。\\ $S_{\triangle POB} = \frac{1}{2} \times OB \times |x_p| = \frac{1}{2} \times 4 \times |x_p| = 2|x_p|$。}
    \solstep{求解未知坐标}{令 $2|x_p| = 6 \implies |x_p| = 3 \implies x_p = 3$ 或 $x_p = -3$。\\ 故点 $P$ 坐标为 $(3, 0)$ 或 $(-3, 0)$。}
  \end{solutionblock}
\end{problem}


\needspace{8\baselineskip}

\begin{problem}[points=20]
  直线 $l_1$ 经过 $A(0, 3)$ 和 $B(2, 0)$ 两点；直线 $l_2$ 经过 $B(2, 0)$ 和 $C(4, 4)$ 两点。

(1) 求直线 $l_1$ 和直线 $l_2$ 的解析式；
(2) 过点 $A$ 作直线 $l_2$ 的平行线交 $x$ 轴于点 $D$，求点 $D$ 的坐标。
  \answerarea[50mm]
  \setcounter{solstep}{0}
  \begin{solutionblock}
    \textbf{答案：}(1) $l_1: y = -\dfrac{3}{2}x + 3$，$l_2: y = 2x - 4$；(2) $D\left(-\dfrac{3}{2}, 0\right)$\par\medskip
    \solstep{求 l1 解析式}{设 $l_1: y = k_1x + b_1$，代入 $A(0,3)$得 $b_1 = 3$。\\ 代入 $B(2,0)$ 得 $2k_1 + 3 = 0 \implies k_1 = -\frac{3}{2}$，故 $l_1: y = -\frac{3}{2}x + 3$。}
    \solstep{求 l2 解析式}{设 $l_2: y = k_2x + b_2$，代入 $B(2,0)$ 和 $C(4,4)$：\\ $\begin{cases} 2k_2 + b_2 = 0 \\ 4k_2 + b_2 = 4 \end{cases} \implies \begin{cases} k_2 = 2 \\ b_2 = -4 \end{cases}$，故 $l_2: y = 2x - 4$。}
    \solstep{求平行线和 D 坐标}{平行于 $l_2$ 且过 $A(0,3)$，其斜率 $k = 2$。\\ 故平行线解析式为 $y = 2x + 3$。\\ 交 $x$ 轴于 $D$：令 $y = 0 \implies 2x + 3 = 0 \implies x = -\frac{3}{2}$。\\ 故点 $D$ 坐标为 $\left(-\frac{3}{2}, 0\right)$。}
  \end{solutionblock}
\end{problem}


\clearpage
\section*{参考答案}


\begin{keyideabox}
  \textbf{选择题}
  1. A  2. A  3. C  4. A
\end{keyideabox}



\begin{keyideabox}
  \textbf{填空题}
  1. 5
\end{keyideabox}



\begin{keyideabox}
  \textbf{解答题}
  第 1 题：(1) $A(3, 0)$，$B(0, 4)$；(2) $P(3, 0)$ 或 $P(-3, 0)$
第 2 题：(1) $l_1: y = -\dfrac{3}{2}x + 3$，$l_2: y = 2x - 4$；(2) $D\left(-\dfrac{3}{2}, 0\right)$
\end{keyideabox}



\end{document}
